% =====================================================================
%  Chapter 11 — Step 5
% =====================================================================

\chapter{Step Five: The Budget, and the Two Thirds}
\label{ch:budget}

This chapter is the centre of the book, because it is where the remark's
most consequential admission is computed --- and, remarkably, the computation is
in the source.

\section{The remark's admission}

The remark's item 14 reads, in full:

\begin{quote}
\small
``precise computation of the `cost' of each lift, showing that one would stop at
2/3 a priori, attempt to get a main term and transfer constant to a function
field or kill via twisting, need for `division'.''
\end{quote}

Read the structure of that sentence carefully, because it is doing four
different things at once and it is easy to read only the first:

\begin{enumerate}
\item A computation was carried out.
\item Its result: the method would stop at two thirds, a priori.
\item The response was an \emph{attempt}, in three named forms.
\item The attempt required a \emph{division}, which is named and not explained.
\end{enumerate}

The remark does not give the computation. It gives the answer. What the
development contains is the computation.

\section{Where the two thirds actually is}

The constant is not a metaphor and it is not the author's summary of a
hand-waving estimate. It is a literal constant in the submitted source, in
\code{Formalism/SharpRestriction/GrowthBounds/SegmentDegree/}, and it occurs in
the threshold arithmetic that decides how much of the available growth the
argument can use.

\begin{seamnote}{The two threshold lemmas (\code{ThresholdPowerBoundOfCounts.lean})}
{\footnotesize
\begin{verbatim}
theorem threshold_arith_letter_case (k Ns c q bo : Real)
(hk5 : 5 <= k) (hN : 1 <= Ns) (hc : 40 * (k * Ns) <= c) (hq
      : 9 / 10 * c <= q)
    (hb : q - k < bo) : c / 2 <= 2 / 3 * bo - k * Ns

theorem threshold_arith_w_case (k Ns c q fl r : Real)
(hk5 : 5 <= k) (hN : 1 <= Ns) (hr : r <= k) (hc : 40 * (k *
      Ns) <= c)
    (hq : 9 / 10 * c <= q) (hfl : q - 1 < fl) :
c / 2 <= 2 / 3 * (fl - (k + r + 1) - (k * Ns + 1)) - k * Ns
\end{verbatim}
}
\end{seamnote}

Let us read the variables.

\begin{description}[style=nextline,leftmargin=0pt,font=\sffamily\bfseries\small,
  itemsep=6pt]
\item[$k$] The weight: $k = r + K$ where $r$ and $K$ are the two rank
parameters of the data. The hypothesis is $k \ge 5$.
\item[$N_s$] The level cost: the size of the integral structure. The
hypothesis is $N_s \ge 1$.
\item[$c$] \textbf{The available growth.} The hypothesis is
$40\,(k N_s) \le c$: the growth must dominate forty times the weight times the
level cost. This is the requirement that the multiplicativity of the chain be
strong enough to pay for the integrality.
\item[$q$] The budget: the hypothesis is $\tfrac{9}{10}c \le q$, so the budget
must be at least nine tenths of the available growth.
\item[$bo$, $fl$] The budget actually reached by the segment, in the two cases.
\end{description}

And the conclusion, in both cases, has the same shape:
\begin{equation}
\frac{c}{2} \;\le\; \frac{2}{3}\cdot(\text{budget reached}) \;-\; k N_s.
\label{eq:threshold}
\end{equation}

Read \eqref{eq:threshold} as the condition the argument must satisfy in order
to conclude. The left-hand side is half the available growth: the amount of
growth the argument needs to dominate in order to get a square-root-quality
remainder. The right-hand side is \emph{two thirds} of the budget the segment
actually reached, minus the weight-times-level-cost that the integrality eats.

\begin{quote}
\textbf{What the two thirds is.} The growth the argument can use is two
thirds of the budget it reaches. The third is consumed. In the remark's own
phrase from item 14, the method ``would stop at 2/3 a priori'' --- the
computation is exactly this inequality, and the number in the remark is this
constant.
\end{quote}

For the second case the budget is itself reduced: the bound obtained is
$\tfrac23\bigl(fl - (k+r+1) - (kN_s+1)\bigr) - kN_s$, so the subtraction of
$k+r+1$ and $kN_s+1$ is charged \emph{before} the two-thirds discount is
applied. That is the "cost of each lift" in the most literal sense available:
each stage of the chain charges a fixed amount in weight and level, and the
discount is applied to what is left.

\section{Where the third goes: the two paths}

The segment-degree machinery explains what the missing third is for. There are
two classes of path through the letters of a segment, and the two thirds
appears as the requirement on both.

\begin{seamnote}{\code{Typ} and \code{GoodA} (from \code{AtypicalPathSmallWeight.lean})}
{\footnotesize
\begin{verbatim}
def Typ [NeZero r] (LL : Letters (Fin N)) (B : Fin N -> Nat)
  (od : Fin N -> Fin N)
(bb : Fin N -> Real) (La : Real) (p : Nat) (mu : Fin (N + 1)
      -> Fin (r + K) -> Int) : Prop :=
  forall i : Fin N,
(LL.vec i = false -> (posS (mu i.castSucc) - posS (mu
      i.succ) : Int) <= Xi r K bb La p i) and
(LL.vec i = true -> r <= i.val -> (negS (mu i.castSucc) -
      negS (mu i.succ) : Int) <= Xi r K bb La p i) and
(r <= i.val -> (2 / 3 : Real) * (B (od i) : Real) <= (plPart
      mu i.castSucc (lastIdx r) : Int))

def GoodA [NeZero r] ... : Prop :=
  forall i : Fin N, r <= i.val ->
(2 / 3 : Real) * (B (od i) : Real) <= (plPart mu i.castSucc
      (lastIdx r) : Int) and
(LL.vec i = true -> (sum c, (tanPart mu i.succ c - tanPart
      mu i.castSucc c) : Int)
        <= Xi r K bb La p i + 2 * (B (od i) : Real)) and
(LL.vec i = false -> (sum c, (plPart mu i.castSucc c -
      plPart mu i.succ c) : Int)
        <= Xi r K bb La p i)
\end{verbatim}
}
\end{seamnote}

Two things to read here.

\begin{description}[style=nextline,leftmargin=0pt,font=\sffamily\bfseries\small,
  itemsep=6pt]
\item[Both paths carry the same two-thirds requirement.]
$(2/3)\cdot B(\mathrm{od}\ i) \le \mathrm{plPart}\ \mu\ (\mathrm{lastIdx}\ r)$:
two thirds of the budget at letter $i$ must be carried by the weight
available at the last index of the segment. This is a constraint on the
\emph{weight}, and the weight is what the chain of lifts is supposed to
supply.
\item[The atypical path gets slack, and it is exactly the slack that is
missing elsewhere.] In \code{GoodA} the second condition carries an extra
$2\cdot B(\mathrm{od}\ i)$: the atypical path is allowed to exceed the
per-letter budget \code{Xi} by twice the budget. The typical path is not.

So the geometry is: the argument proceeds by assuming the segment follows a
typical path, where each letter's contribution is tightly controlled. The
atypical path --- the one where a letter misbehaves --- must be handled
separately, and handling it costs. The cost is paid out of the budget, and
what survives is two thirds.
\end{description}

The supporting files are named for exactly this: \code{GoodPathWeight.lean},
\code{TypOfTypical.lean}, \code{TypicalPathOrdersSmall.lean},
\code{AtypicalPathSmallWeight.lean}, \code{TopOccursInAtMostTwoSegments.lean},
\code{WindowSmall.lean}, \code{AllowedLettersOfLateRows.lean}.

\section{What would have to change to cross it}

Now the arithmetic of crossing, which the remark's three remedies map onto
precisely.

The condition to be proved is \eqref{eq:threshold}. To get the full budget
rather than two thirds of it, one needs
\begin{equation}
\frac{c}{2} \;\le\; (\text{budget reached}) \;-\; k N_s
\label{eq:thresholdfull}
\end{equation}
instead of \eqref{eq:threshold}. The difference between the two right-hand
sides is exactly $\tfrac13$ of the budget reached. So the method needs to find
one third more of the budget, and the remark's three attempts are three ways
of trying to find it.

\begin{description}[style=nextline,leftmargin=0pt,font=\sffamily\bfseries\small,
  itemsep=6pt]
\item[``Get a main term.'] More of the available growth $c$. The hypothesis
$40(kN_s) \le c$ says the growth must dominate the weight and level; getting
more growth means either more lifts or a better per-lift multiplicativity.
\item[``Transfer constant to a function field.'] Move the constant from the
integer setting to a function field, where the Weil machinery is native and
the arithmetic factor exists. This is the same move as the function-field
proof of the Weil bound, applied to the constant rather than to the whole
argument.
\item[``Kill via twisting.'] Eliminate the atypical path by twisting, so that
the third is never spent. This is the same device used in
Chapter~\ref{ch:vanishing} to kill cohomology, applied earlier.
\item[``Division.''] And the division: something has to be divided, and the
development has a file named for it,
\code{Formalism/SharpRestriction/GrowthBounds/SegmentDegree/DivMulSubLeDivSubDiv.lean}.
\end{description}

The remark's word ``division'' is the one place where the summary is most
opaque and the source is most concrete, and it is worth a reader's attention
precisely because it is the step that would have to work for the claim to be
true.

\section{The two-thirds reading against the August theorem}

There is one more thing to check, and it is the check that distinguishes a
coincidence from a mechanism.

The remark's author had already published a two-thirds theorem three months
earlier --- arXiv:2608.13637, ``More than two thirds of the zeta zeros are
simple and on the critical line'' --- using the same method family. The
fraction 2/3 in the growth bounds and the fraction 2/3 in the published
theorem are the same number. Chapter~\ref{ch:twothirds} treats that theorem.

\begin{quote}
The natural reading is that the budget computation of item 14 is the
explanation of the two-thirds theorem: the method, costed, delivers a bound
covering two thirds of the zeros, which is precisely what the August paper
proves, and the three attempts are what was done about the remaining third.
\end{quote}

This is a coherent, checkable story, and it explains why the author would
compute a ceiling before announcing a full result. It does not, by itself,
establish that the remaining third was crossed. It is the central question the
book keeps coming back to, and Chapter~\ref{ch:verdict} states it as such.

\begin{keytakeaways}
\item The two thirds is a literal constant in the submitted source, in the
segment-degree threshold arithmetic that decides how much of the available
growth the argument can use.
\item The threshold condition is $c/2 \le \tfrac23\cdot(\text{budget reached})
- kN_s$, where $c$ is the available growth, the budget is what the segment
reaches, and $kN_s$ is the weight-and-level cost charged at each stage.
\item The third is consumed by the atypical path: the typical path controls
each letter tightly, the atypical path is allowed an extra $2B$ of slack, and
what survives is two thirds of the budget.
\item Crossing the ceiling means finding one third more of the budget --- which
is what the remark's three attempts, and its ``division'', are for.
\item The same fraction appears in the author's own published two-thirds
theorem of three months earlier, which is the strongest available indication
that the ceiling is real and that item 14 is its explanation.
\end{keytakeaways}
